\begin{examplebox}
	\textbf{\textcolor{CExample}{例 1（基础）}}

	\textbf{\textcolor{CExample}{题目：}}
	已知函数$f(x)$满足$f(x+2)=\frac{1}{f(x)}$（$f(x)\neq 0$），且$f(1)=2$，求$f(5)$的值。

	\textbf{\textcolor{CExample}{解题步骤：}}
	\begin{description}[style=nextline, leftmargin=2.8em, labelsep=0.5em, itemsep=0.45em]
		\item[\textbf{第一步（识别条件）：}] 条件为$f(x+2)=\frac{1}{f(x)}$，其中$a=2$，$f(1)=2\neq0$，符合倒数型周期条件。
			\par\textbf{理由：} 匹配结论条件$f(x+a)=1/f(x)$，$a=2$。
		\item[\textbf{第二步（求周期）：}] 由结论得周期$T=2a=4$。
			\par\textbf{理由：} 结论直接给出。
		\item[\textbf{第三步（周期转化）：}] $5=1+4$，所以$f(5)=f(1)$。
			\par\textbf{理由：} 周期为4，自变量相差4函数值相等。
		\item[\textbf{第四步（代入即得）：}] $f(5)=f(1)=2$。
			\par\textbf{理由：} 已知$f(1)=2$。
	\end{description}

	\textbf{\textcolor{CExample}{关键结论：}}
	\textbf{$f(5)=2$。}

	\medskip
	\textbf{\textcolor{CExample}{例 2（稍变形）}}

	\textbf{\textcolor{CExample}{题目：}}
	已知函数$f(x)$满足$f(x+3)=-\frac{1}{f(x)}$（$f(x)\neq 0$），且$f(0)=5$，求$f(6)$的值。

	\textbf{\textcolor{CExample}{解题步骤：}}
	\begin{description}[style=nextline, leftmargin=2.8em, labelsep=0.5em, itemsep=0.45em]
		\item[\textbf{第一步（识别条件）：}] 条件为$f(x+3)=-\frac{1}{f(x)}$，属于负倒数型，$a=3$。
			\par\textbf{理由：} 匹配结论条件$f(x+a)=-1/f(x)$。
		\item[\textbf{第二步（求周期）：}] 由结论得周期$T=2a=6$。
			\par\textbf{理由：} 结论直接适用。
		\item[\textbf{第三步（周期转化）：}] $6=0+6$，故$f(6)=f(0)$。
			\par\textbf{理由：} 周期性。
		\item[\textbf{第四步（代入即得）：}] $f(6)=f(0)=5$。
			\par\textbf{理由：} 已知$f(0)=5$。
	\end{description}

	\textbf{\textcolor{CExample}{关键结论：}}
	\textbf{$f(6)=5$。}

	\medskip
	\textbf{\textcolor{CExample}{例 3（综合应用）}}

	\textbf{\textcolor{CExample}{题目：}}
	已知函数$f(x)$满足$f(x+1)=\frac{1}{f(x)}$（$f(x)\neq 0$），且当$x\in[0,1]$时$f(x)=2^x$，求$f(2023.5)$的值。

	\textbf{\textcolor{CExample}{解题步骤：}}
	\begin{description}[style=nextline, leftmargin=2.8em, labelsep=0.5em, itemsep=0.45em]
		\item[\textbf{第一步（求周期）：}] 由$f(x+1)=\frac{1}{f(x)}$，得$a=1$，周期$T=2a=2$。
			\par\textbf{理由：} 直接套用结论。
		\item[\textbf{第二步（周期转化）：}] $2023.5=2\times1011+1.5$，故$f(2023.5)=f(1.5)$。
			\par\textbf{理由：} 周期函数性质：$f(x+2k)=f(x)$。
		\item[\textbf{第三步（原关系转化）：}] 利用条件$f(x+1)=\frac{1}{f(x)}$，取$x=0.5$得$f(1.5)=\frac{1}{f(0.5)}$。
			\par\textbf{理由：} 条件对任意$x$成立。
		\item[\textbf{第四步（区间求值）：}] 因为$0.5\in[0,1]$，由已知$f(0.5)=2^{0.5}$，即 $\sqrt{2}$。
			\par\textbf{理由：} 分段表达式。
		\item[\textbf{第五步（计算结果）：}] $f(2023.5)=\frac{1}{\sqrt{2}}$，即 $\frac{\sqrt{2}}{2}$。
			\par\textbf{理由：} 代入化简。
	\end{description}

	\textbf{\textcolor{CExample}{关键结论：}}
	\textbf{$f(2023.5)=\dfrac{\sqrt{2}}{2}$。}
\end{examplebox}
